% Part: many-valued-logic
% Chapter: sequent-calculus
% Section: introduction

\documentclass[../../../include/open-logic-section]{subfiles}

\begin{document}

\olfileid[hi]{mvl}{seq}{int}

\olsection{परिचय}

शास्त्रीय तर्कशास्त्र का सीक्वेंट कलन एक दक्ष और सरल !!{derivation}-तंत्र
है। यदि कोई बहुमूल्य तर्कशास्त्र परिमित संख्या में सत्य-मूल्यों वाले आव्यूह
से परिभाषित हो, अर्थात् $V$ परिमित हो, तो उसके लिए सीक्वेंट कलन दिया जा
सकता है। इसे करने का विचार शास्त्रीय मामले में सीक्वेंटों के अर्थ और
अनुमान-नियमों के रूप पर विचार करने से मिलता है।

अब स्मरण करें कि एक सीक्वेंट
\begin{align*}
    {}!A_1, \dots, !A_n & \Sequent !B_1, \dots, !B_n
\intertext{की व्याख्या इस !!{formula} के रूप में की जा सकती है:}
    (!A_1 \land \cdots \land !A_m) & \lif (!B_1 \lor \cdots \lor
{}!B_n).
\end{align*}
दूसरे शब्दों में, कोई !!{valuation}~$\pAssign{v}$ किसी सीक्वेंट
$\Gamma \Sequent \Delta$ को तभी \emph{संतुष्ट करता है} जब या तो किसी
$!A \in \Gamma$ के लिए $\pValue{v}(!A) = \False$ हो, या किसी
$!A \in \Delta$ के लिए $\pValue{v}(!A) = \True$ हो। इस व्याख्या पर आरंभिक
सीक्वेंट $!A \Sequent !A$ सदैव संतुष्ट होते हैं, क्योंकि या तो
$\pValue v(!A) = \True$ या $\pValue(!A) = \False$।

यहाँ पार्श्व सूत्रों $\Gamma$, $\Delta$ को हटाकर $\Log{LK}$ में सशर्त के
अनुमान-नियम दिए गए हैं:

\begin{defish}
    \Axiom$ \fCenter !A$
    \Axiom$ !B \fCenter $
    \RightLabel{\LeftR{\lif}}
    \BinaryInf$ !A \lif !B \fCenter $
    \DisplayProof
    \hfill
    \Axiom$ !A \fCenter !B$
    \RightLabel{\RightR{\lif}}
    \UnaryInf$ \fCenter !A \lif !B $
    \DisplayProof
\end{defish}

यदि हम ऊपर दी गई सीक्वेंट की आर्थी व्याख्या लागू करें, तो
$\LeftR{\lif}$ नियम कहता है कि यदि $\pValue v(!A) = \True$ और
$\pValue v(!B) = \False$, तो $\pValue v(!A \lif !B) = \False$। इसी तरह
$\RightR{\lif}$ नियम कहता है कि यदि या तो $\pValue v(!A) = \False$ या
$\pValue v(!B) = \True$, तो $\pValue v(!A \lif !B) = \True$। और वास्तव
में ये सशर्त कथन द्विसशर्त कथन हैं। $\LeftR{\land}$ और $\RightR{\lor}$
नियमों के मानक सूत्रीकरण में संगत सशर्त कथन द्विसशर्त नहीं होंगे। किंतु इन
नियमों के ऐसे वैकल्पिक रूप हैं जिनमें वे द्विसशर्त होते हैं:

\begin{defish}
    \Axiom$!A, !B, \Gamma \fCenter \Delta$
    \RightLabel{\LeftR{\land}}
    \UnaryInf$!A \land !B, \Gamma \fCenter \Delta$
    \DisplayProof
    \hfill
    \Axiom$ \Gamma \fCenter \Delta, !A, !B$
    \RightLabel{\RightR{\lor}}
    \UnaryInf$ \Gamma \fCenter \Delta, !A \lor !B$
    \DisplayProof
\end{defish}

जब इस मूल विचार को किसी $n$-मूल्य तर्कशास्त्र पर लागू किया जाता है, तो दो
के स्थान पर $n$ स्थानों वाला सीक्वेंट कलन प्राप्त होता है---प्रत्येक
सत्य-मूल्य के लिए एक स्थान। $V = \{\False, \Undef, \True\}$ वाले किसी
त्रिमूल्य तर्कशास्त्र में सीक्वेंट $\Gamma \mid \Pi \mid \Delta$ रूप का
व्यंजक होता है। कोई !!{valuation}~$\pAssign v$ इसे तभी संतुष्ट करता है जब या
तो किसी $!A \in \Gamma$ के लिए $\pValue{v}(!A) = \False$, या किसी
$!A \in \Delta$ के लिए $\pValue{v}(!A) = \True$, या किसी $!A \in \Pi$
के लिए $\pValue{v}(!A) = \Undef$ हो। फलतः आरंभिक सीक्वेंट
$!A \mid !A \mid !A$ सदैव संतुष्ट होते हैं।

\end{document}
